% Fragmento autónomo para el borrador del artículo II, 8 de septiembre de 2026.
% Procedencia: RESULTADO_RECUPERADO y FORMALIZACION_REUNIDA de los propietarios
% B07/c41_ckm.tex y C02/74_mezcla_contenido_exclusivo_post_rev2.tex del
% BORRADOR_02; integral 2249, caps. 49, 92.6, 92.24 y 93.10--93.11.
% Los originales y su soporte demostrativo íntegro no se sustituyen ni editan.
% Claves bibliográficas utilizadas: ckm, mezcla, pmnsespectral, cpt.
% Necesita: amsmath, amssymb, amsthm; theorem/proposition/definition/remark.

\section{Angles de mélange CKM, phase CP et transport leptonique}
\label{sec:ii-propagacion-angular}

\subsection{Antécédent commun et séparation des deux transports}

Le fondement commun de l’article produit l’état enrichi au moyen
des deux évaluations d’APP, de la sélection orientée du TRIT et de la
composition effective de sélection, transport, retenue, mémoire et retour
du TPK. Les représentations angulaires ne remplacent pas cet état : ce sont des lecteurs
de son registre. En particulier, le prolongement
\(w_6\to w_{12}\to w_{18}\to w_{24}\to w_{30}\to R_{36}\to G_9\)
conserve la provenance des lectures numériques, et l’holonomie
\(\Gamma_9\) agit avant la représentation réduite \(K_{\rm ph}\).
Le retour de phase ne réinitialise pas le registre.

De ce même antécédent, on reçoit ici deux classes d’objets. La première
contient l’angle direct, le contre-angle et le défaut torsionnel
\(\Delta_4\), déjà construits dans leurs lecteurs propres. La seconde contient
les extensions survivantes, leurs chemins observables, le registre
dodécaphasique et les opérateurs des fibres matérielles. La chaîne
extension--chemin--registre--signature--caractère--tours conserve l’atlas
\(81/56/13/324\) et distingue la coordinatisation nulle de la branche positive.
Les angles de mélange ne choisissent pas rétrospectivement ces chemins et n’engendrent pas
les valeurs propres de masse.

Le transport des quarks utilise une coordinatisation angulaire rationnelle et sa réalisation
unitaire complexe. Le transport leptonique utilise les repères spectraux
des fibres chargée et neutre. Tous deux proviennent de la structure discrète
commune, mais leurs domaines, leurs lecteurs et leurs changements de base sont distincts.
La fonctionnelle de Barbero--Immirzi et le mélange ne se définissent pas non plus l’un par
l’autre : ils reçoivent des représentations typées de l’antécédent angulaire et d’incidence.
Leur relation s’exprime au moyen de ces applications, non au moyen d’une
formule rétrospective utilisant une valeur de Barbero comme sélecteur de mélange.

\begin{remark}[Règles sectorielles et évaluation du lecteur]
Les trois règles sectorielles et leur évaluation sont développées ci-après
dans la section~\ref{sec:ii09-construccion-sectorial}. Le théorème
\ref{thm:ii09-matriz-sectorial} obtient les colonnes comme images de la
base angulaire. Les preuves ultérieures établissent l’inversion de la coordinatisation,
le transport de feuillet, la métrique et la réalisation complexe.
Le choix des règles et l’évaluation de leurs coefficients sont deux
questions mathématiques distinctes : le système de départ est déclaré dans la
Définition~\ref{def:ii09-sistema-sectorial}, et la preuve du théorème
utilise ce système complet. La composition préalable qui sélectionne ces
règles à partir de l’incidence nécessite sa démonstration propre.
\end{remark}

% Aportación separada, 9 de septiembre de 2026.
% RESULTADO_RECUPERADO: construcción sectorial del propietario 08_ckm_cp.tex.
% FORMALIZACION_REUNIDA: descomposición de paridad y compatibilidad representacional.
% Insertar antes de la carta racional CKM; conservar el desarrollo posterior.
% La procedencia y el alcance exacto están en 08b_procedencia_angular.md.

\subsection{Incidence sectorielle et construction du lecteur angulaire}
\label{sec:ii09-construccion-sectorial}

La construction angulaire utilise les représentations d’action et de clôture
obtenues au préalable à partir d’APP--TRIT--TPK. L’angle direct, le
contre-angle et le défaut torsionnel conservent leurs domaines propres.
Le lecteur des quarks combine ces trois coordonnées au moyen de l’incidence
sectorielle ; son évaluation précède les amplitudes complexes et la
comparaison métrologique. L’inclusion hexadique--octadique intervient dans
le troisième secteur à travers la différence de leurs cardinaux.

% ARQUITECTURA_AUTORAL_PREEXISTENTE: semicuadrado bidireccional de N42.
% FORMALIZACION_NUEVA: grupoide de pares y ponderacion por estabilizadores.
\subsubsection{Incidence marquée et normalisation bidirectionnelle}
\label{sec:ii09-normalizacion-bidireccional}

Les invariants du lecteur conservent la nature de l'objet dont ils
proviennent. L'intersection marquée $B_K=H_e\cap P_\pi$ a pour cardinal
$b=4$ ; l'entier $q=7$ est la coordonnée du pivot $p_\varphi$ dans la
numérotation orientée. L'ordre de détermination du système de Steiner
est $p=5$. Les supports hexadique et octadique ont pour cardinaux $h_6=6$
et $o_8=8$, et l'alphabet d'orientation tritique a pour cardinal $t=3$.
Un cardinal de support et une coordonnée de pivot sont utilisés par
leurs applications respectives, même lorsqu'ils sont tous deux publiés comme entiers.

La normalisation de l'autocroisement pentagonal admet la
réalisation explicite suivante. Soit $P$ le support de cinq positions de la clôture
et soit $X=P\times P$ son espace de paires ordonnées. L'inversion
bidirectionnelle sur cet espace est
\[
 \jmath:X\longrightarrow X,\qquad \jmath(x,y)=(y,x),
 \qquad\jmath^2=I.
\]
L'information de stabilisateur de chaque orbite est conservée. Pour une
orbite $[z]$, on définit son poids par
$1/|\operatorname{Stab}_{C_2}(z)|$. La somme de ces poids est la
cardinalité du groupoïde d'action $X/\!/C_2$.

\begin{proposition}[Mesure orbitale de l'autocroisement bidirectionnel]
\label{prop:ii09-semicuadrado-orbital}
Pour un support fini de cardinal $p$, l'inversion des paires ordonnées
détermine
\[
 |(P\times P)/\!/C_2|
 =\sum_{[z]\in(P\times P)/C_2}
      \frac1{|\operatorname{Stab}_{C_2}(z)|}
 =\frac{p^2}{2}.
\]
Pour $p=5$, l'autocroisement possède dix orbites libres et cinq orbites
fixes, et sa mesure orbitale est $10+5/2=25/2$.
\end{proposition}
\begin{proof}
Chaque paire diagonale $(x,x)$ forme une orbite fixe, dont le stabilisateur
est d'ordre deux ; il y en a $p$. Les $p(p-1)$ paires restantes se
regroupent deux par deux, avec un stabilisateur trivial. Leur contribution est
donc $p(p-1)/2$, et celle de la diagonale est $p/2$. La somme vaut
$p^2/2$. De manière équivalente, l'identité orbite--stabilisateur donne
$1/|\operatorname{Stab}(z)|=|[z]|/2$ ; la sommation sur toutes les orbites
produit $|P\times P|/2$.
\end{proof}

La mesure conserve explicitement la contribution des points fixes.
L'ensemble ordinaire des orbites a pour cardinal $p(p+1)/2$, qui vaut
quinze dans la clôture pentagonale. La mesure orbitale et ce cardinal sont
deux lectures déterminées de la même inversion. La normalisation
bidirectionnelle du lecteur utilise la première : les dix orbites libres
pèsent une unité et les cinq diagonales pèsent une demi-unité.

Cette réalisation précise le facteur $p^2/2$ employé dans le système
sectoriel suivant. L'affectation de cette contribution au secteur $23$,
son orientation négative et les deux autres règles d'incidence restent
spécifiées par le système sectoriel complet. La démonstration établit la
normalisation de l'autocroisement et maintient distincte la sélection des
fonctionnelles angulaires.


% FORMALIZACION_NUEVA de la elevacion de N39, compatible con (rho9,q9),
% T_d y Upd_t; arquitectura de residuo/cociente y orientacion preexistente.
% Fragmento de investigacion separado. No selecciona las reglas CKM N42.
\subsubsection{Transport de retenue arithmétique et holonomie entière du croisement arithmétique}
\label{subsec:transporte-entero-cruce}

La lecture modulaire du croisement admet une réalisation entière qui maintient
l’information accumulée lors du franchissement du bord. On développe ici cette
réalisation au moyen des quotients de transport. Son domaine est une coordinatisation
relevée du support APP ; l’assignation ultérieure d’une région à un secteur
angulaire conserve sa propre application d’incidence.

\paragraph{Coordinatisation nonadique et transport relevé.}
On emploie la décomposition exacte
\[
 \rho_9(n)=1+((n-1)\bmod9),\qquad
 q_9(n)=\left\lfloor\frac{n-1}{9}\right\rfloor,
 \qquad n=\rho_9(n)+9q_9(n).
\]
Une position relevée est
$z=(x,y,q_x,q_y)\in D_9^2\times\mathbb Z^2$, où
$X=x+9q_x$ et $Y=y+9q_y$. L'incrément entier $u$ de la première
coordonnée agit par
\[
 T^x_u(x,y,q_x,q_y)=
 \bigl(\rho_9(x+u),y,q_x+c_9(x;u),q_y\bigr),\qquad
 c_9(x;u)=\left\lfloor\frac{x+u-1}{9}\right\rfloor.
\]
La deuxième coordonnée se transporte de la même manière. Les avances
cardinales correspondent à $u=\pm1$ ; l'incrément nul conserve l'état.
Le cocycle satisfait
\[
 c_9(x;u+v)=c_9(x;u)+c_9(\rho_9(x+u);v).
\]
En effet, les deux membres sont le quotient de l'unique décomposition de
$x+u+v$ avec résidu dans $D_9$. Par conséquent,
$T^x_vT^x_u=T^x_{u+v}$ et $T^x_{-u}=(T^x_u)^{-1}$ ; les transports
des deux coordonnées commutent. L'état conserve le nombre de franchissements
du bord, même lorsque la position visible revient.

\paragraph{Feuillets relevés et mise à jour des registres.}
La réalisation entière est fixée en prolongeant les opérations arithmétiques
APP aux coordonnées transportées :
\begin{align*}
 \widetilde S(z)&=X+Y=x+y+9(q_x+q_y),\\
 \widetilde P(z)&=XY
 =xy+9(q_xy+q_yx)+81q_xq_y.
\end{align*}
À chaque position, le registre des deux feuillets est conservé :
\[
 \bigl(\rho_9(\widetilde S),q_9(\widetilde S),
       \rho_9(\widetilde P),q_9(\widetilde P)\bigr).
\]
Si $e:z\to z'$ est un lien et $T$ le feuillet sélectionné, son incrément
intégral s'obtient à partir du registre avant et après le transport :
\begin{equation}
 \delta_e\widetilde T=
 \rho_9(\widetilde T(z'))-\rho_9(\widetilde T(z))
 +9\bigl(q_9(\widetilde T(z'))-q_9(\widetilde T(z))\bigr).
 \label{eq:cruce-incremento-registro}
\end{equation}
Cette identité explicite la conservation : le saut du résidu et
l’accroissement du quotient reconstruisent ensemble la variation du feuillet.
La sélection de feuillet, le transport cardinal et la mise à jour du
registre sont des opérations distinctes. Dans cette coordinatisation, on compose d’abord la
sélection ordonnée de feuillets $(T_x,T_y)$ avec $T_d$ et avec la mise à jour
de leurs quotients ; la formule ne postule pas que tout le TPK soit réduit
à cette représentation locale.

\begin{proposition}[Courbure entière et transport de frontière]
\label{prop:cruce-stokes-integral}
Définissons
$\widetilde A_x=\delta_x\widetilde T_x$ et
$\widetilde A_y=\delta_y\widetilde T_y$, au moyen de
\eqref{eq:cruce-incremento-registro}, et soit
\[
 \widetilde F_{xy}(X,Y)=
 \widetilde A_x(X,Y)+\widetilde A_y(X+1,Y)
 -\widetilde A_x(X,Y+1)-\widetilde A_y(X,Y).
\]
Les choix homogènes ont une courbure nulle. Les choix mixtes
$(\widetilde S,\widetilde P)$ et
$(\widetilde P,\widetilde S)$ ont respectivement pour courbures les entiers
$+1$ et $-1$. Pour le bord positif d'un rectangle
$R_{a,b}$ de côtés entiers $a,b>0$,
\[
 \widetilde W_{SP}(\partial R_{a,b})=ab,\qquad
 \widetilde W_{PS}(\partial R_{a,b})=-ab.
\]
Le résultat est indépendant du point de base et du nombre de franchissements
des raccords nonadiques.
\end{proposition}

\begin{proof}
La reconstruction exacte des deux feuillets donne
\[
 \delta_x\widetilde S=\delta_y\widetilde S=1,
 \qquad\delta_x\widetilde P=Y,
 \qquad\delta_y\widetilde P=X.
\]
Ainsi, la connexion $SP$ est $(1,X)$, de courbure
$1+(X+1)-1-X=1$ ; la connexion $PS$ est $(Y,1)$, de courbure
$Y+1-(Y+1)-1=-1$. Pour un feuillet unique, les différences
mixtes s'annulent. La somme des courbures des $ab$ plaquettes annule chaque
lien intérieur avec son inverse et conserve les liens de frontière.
On obtient ainsi $\pm ab$ dans $\mathbb Z$. Les quotients transportés
garantissent que la même identité vaut lors du franchissement de tout raccord.
\end{proof}

L'accumulateur de frontière peut à son tour être conservé sous forme nonadique :
$W=r_W+9q_W$. Pour une contribution entière $w_e$, il est mis à jour par
\[
 r_W'=\rho_9(r_W+w_e),\qquad
 q_W'=q_W+\left\lfloor\frac{r_W+w_e-1}{9}\right\rfloor.
\]
La composition par liens fournit la même somme que son évaluation
entière directe. Dans cette coordinatisation, le zéro entier a pour coordonnées
$(r_W,q_W)=(9,-1)$ ; c’est une représentation arithmétique de l’accumulateur,
sans identification supplémentaire avec les différents zéros orientés du TRIT.

\paragraph{Trois opérations sur l'orientation.}
Soit $\gamma$ une route cardinale fermée. Son inversion temporelle
$\gamma^{-1}$ inverse l'ordre des liens et l'orientation de chacun ;
par additivité,
\[
 \widetilde W_{SP}(\gamma^{-1})=-\widetilde W_{SP}(\gamma).
\]
L'échange des feuillets satisfait également
\[
 \widetilde W_{PS}(\gamma)=-\widetilde W_{SP}(\gamma).
\]
Pour le vérifier sur toute route, observons que la somme des deux
connexions est l'incrément exact de
$\widetilde S+\widetilde P=X+Y+XY$ ; son intégrale sur une boucle est nulle.

Le demi-tour simultané des deux directions agit, d’autre part,
par $J_c(X,Y)=(2c_x-X,2c_y-Y)$, avec $2c\in\mathbb Z^2$,
de sorte qu’il préserve la coordinatisation du réseau. Pour un lacet, on a
\[
 \widetilde W_{SP}(J_c\gamma)=\widetilde W_{SP}(\gamma).
\]
En effet, la contribution horizontale fermée est nulle et la verticale
se transforme comme
$\sum (2c_x-X)(-\delta Y)=\sum X\delta Y$, car
$\sum\delta Y=0$. Un demi-tour spatial préserve donc
l’orientation aréale de cette coordinatisation. Sa combinaison avec l’inversion temporelle
ou avec l’échange de feuillets se calcule en composant les opérations
correspondantes.

La différence est pertinente pour le registre électronique : l'inversion
simultanée des deux déplacements tous les $27$ pas d'avancement produit, dans
des fenêtres de $9$, la suite
\[
 \tau_j=(-1)^{\lfloor j/3\rfloor},\qquad 0\leq j<12,
 \qquad (+,+,+,-,-,-,+,+,+,-,-,-).
\]
Cette suite enregistre les orientations d'avancement. L'holonomie requiert
en outre les liens parcourus, les feuillets sélectionnés et les quotients
accumulés. Ainsi, le transport de l'orientation électronique est conservé
sans identifier un demi-tour des deux axes à l'inversion temporelle
d'une boucle.

\paragraph{Application au croisement marqué et portée sectorielle.}
Pour le croisement rectangulaire de côtés $4$ et $7$, cette réalisation évalue
exactement
\[
 \widetilde W_{SP}=28=1+9\cdot3,\qquad
 \widetilde W_{SP}(\gamma^{-1})=-28=8+9(-4).
\]
L’entier conserve l’aire orientée ; le résidu isolé représente seulement
sa classe nonadique. Par exemple, à partir de $(7,1)$, les représentants
périodiques des liens ont pour somme $-35$, tandis que la correction de
retenue ajoute $63$ et retrouve $28$.

Dans l'incidence marquée du lecteur angulaire, le $4$ est
$|H_e\cap P_\pi|$ et le $7$ est le pivot marqué $p_\varphi$.
La lecture rectangulaire fournit une réalisation entière du produit
$|H_e\cap P_\pi|\,p_\varphi$ une fois ce croisement déclaré.
L'application qui identifie cette région au secteur $12$ doit conserver
la marque et l'orientation du système sectoriel de
la Définition~\ref{def:ii09-sistema-sectorial}. Les affectations
sectorielles $12,23,13$ et leurs combinaisons de $A,h,\Delta^\circ$
sont une spécification supplémentaire à la loi d'holonomie rectangulaire.
Le présent résultat matérialise le relèvement d'aire, la retenue arithmétique et
leurs involutions, avec ces données de départ explicites.


\begin{definition}[Système sectoriel d'incidence angulaire]
\label{def:ii09-sistema-sectorial}
Le système sectoriel se compose du sextuplet ordonné d'invariants
\[
 (b,p,q,h_6,o_8,t)=(4,5,7,6,8,3)
\]
et des applications linéaires, sur les coordonnées angulaires
\((a,u,d)\),
\begin{align}
 \mathscr A_{12}^{\rm ang}(a,u,d)&=2a-pu+bq\,d,\label{eq:ii09-sector12}\\
 \mathscr A_{23}^{\rm ang}(a,u,d)&=qu-\frac{p^2}{2}\,d,\label{eq:ii09-sector23}\\
 \mathscr A_{13}^{\rm ang}(a,u,d)&=a-pb\,u-\frac{o_8-h_6}{t}\,d.
 \label{eq:ii09-sector13}
\end{align}
Les symboles \(h_6,o_8\) conservent les cardinaux de l'incidence
marquée ; \(t=3\) est la normalisation tritique déclarée de ce secteur.
Le sextuplet et les trois règles forment conjointement le système de
départ de cette construction.
\end{definition}

\begin{theorem}[Évaluation exacte du système sectoriel]
\label{thm:ii09-matriz-sectorial}
L'application conjointe
\((\mathscr A_{12}^{\rm ang},\mathscr A_{23}^{\rm ang},
\mathscr A_{13}^{\rm ang})^{\mathsf T}\) de la
Définition~\ref{def:ii09-sistema-sectorial} a pour unique matrice
\[
 \mathsf M_{\rm CKM}=
 \begin{pmatrix}
 2&-5&28\\
 0&7&-25/2\\
 1&-20&-2/3
 \end{pmatrix}.
\]
Les trois colonnes sont les images des vecteurs de coordonnées
\((1,0,0)\), \((0,1,0)\) et \((0,0,1)\) par ces applications.
\end{theorem}
\begin{proof}
Les opérations sectorielles donnent
\[
 bq=4\cdot7=28,\qquad p^2/2=5^2/2=25/2,
 \qquad pb=5\cdot4=20,\qquad
 (o_8-h_6)/t=(8-6)/3=2/3.
\]
Les images successives de la base de coordonnées sont donc
\[
 c_A=(2,0,1)^{\mathsf T},\qquad
 c_h=(-5,7,-20)^{\mathsf T},\qquad
 c_\Delta=(28,-25/2,-2/3)^{\mathsf T}.
\]
Une application linéaire est déterminée par ses images sur une
base. La matrice dont les colonnes sont ces images est celle indiquée.
La démonstration utilise le système sectoriel complet de la définition.
\end{proof}

Sur les coordonnées précédemment produites, on pose
\[
 (a,u,d)=\left(A_{\deg},\frac{C^*_{\deg}}6,
                       \frac{180}{\pi}\Delta_4\right).
\]
Les contributions directe, conjuguée et torsionnelle à chaque angle
sont ainsi individualisées. En particulier, la différence
\(o_8-h_6=2\) est conservée dans le coefficient du défaut torsionnel
du secteur \(13\). Cette présence concrète permet de suivre
l'incidence jusqu'au lecteur angulaire, en maintenant distincts le
système sectoriel, son évaluation et la paramétrisation unitaire.


\subsection{Coordinatisation rationnelle CKM et unités angulaires}

Notons \(A_{\rm deg}\) et \(C^*_{\rm deg}\) les coordonnées en
degrés des angles HMT déjà engendrés. On utilise les abréviations
\begin{equation}
 h_\theta=\frac{C^*_{\rm deg}}6,
 \qquad
 \Delta_{\rm deg}=\frac{180}{\pi_{\rm HMT}}\Delta_4,
 \qquad
 x_\theta=(A_{\rm deg},h_\theta,\Delta_{\rm deg})^{\mathsf T}.
 \label{eq:ii-angular-chart}
\end{equation}
Ici, \(h_\theta\) est une coordonnée angulaire, distincte de la constante de
Planck \(h\). La conversion en degrés est un changement de coordonnées ultérieur ; elle
n’introduit pas de valeur conventionnelle de \(\pi\) dans le générateur.

\begin{definition}[Lecteur angulaire des quarks]
Le lecteur hérité \(\mathsf M_{\rm CKM}:\mathbb R^3\to\mathbb R^3\)
évalue les coordonnées
\begin{equation}
 \boldsymbol\theta_{\rm deg}=\mathsf M_{\rm CKM}x_\theta,
 \qquad
 \mathsf M_{\rm CKM}=
 \begin{pmatrix}
 2&-5&28\\
 0&7&-25/2\\
 1&-20&-2/3
 \end{pmatrix},
 \qquad
 \delta_{\rm deg}=9A_{\rm deg}.
 \label{eq:ii-ckm-map}
\end{equation}
L'ordre des composantes est \((12,23,13)\). Les fonctions
trigonométriques ultérieures sont évaluées en
\(\vartheta_{ab}=(\pi_{\rm HMT}/180)\theta_{ab,\rm deg}\) et
\(\delta=(\pi_{\rm HMT}/180)\delta_{\rm deg}\).
\end{definition}

\begin{proposition}[Reconstruction rationnelle des coordonnées]
Le lecteur a pour déterminant \(-3857/6\) et pour inverse
\begin{equation}
 \mathsf M_{\rm CKM}^{-1}=
 \begin{pmatrix}
 1528/3857&3380/3857&801/3857\\
 75/3857&176/3857&-150/3857\\
 6/551&-30/551&-12/551
 \end{pmatrix}.
 \label{eq:ii-ckm-inverse}
\end{equation}
En particulier,
\begin{align}
 A_{\rm deg}
 &=\frac{1528\theta_{12,\rm deg}+3380\theta_{23,\rm deg}
          +801\theta_{13,\rm deg}}{3857},\nonumber\\
 C^*_{\rm deg}
 &=\frac{6(75\theta_{12,\rm deg}+176\theta_{23,\rm deg}
          -150\theta_{13,\rm deg})}{3857},\nonumber\\
 \Delta_{\rm deg}
 &=\frac{6\theta_{12,\rm deg}-30\theta_{23,\rm deg}
          -12\theta_{13,\rm deg}}{551}.
 \label{eq:ii-ckm-recover}
\end{align}
La phase satisfait le contrôle linéaire exact
\begin{equation}
 3857\delta_{\rm deg}
 =9(1528\theta_{12,\rm deg}+3380\theta_{23,\rm deg}
      +801\theta_{13,\rm deg}).
 \label{eq:ii-ckm-phase-check}
\end{equation}
\end{proposition}

\begin{proof}
Le développement du déterminant de \eqref{eq:ii-ckm-map} donne
\(-3857/6\). La multiplication de la matrice indiquée dans
\eqref{eq:ii-ckm-inverse} par \(\mathsf M_{\rm CKM}\), dans les deux
ordres, donne \(I_3\). Ses lignes restituent les trois composantes de
\(x_\theta\) ; la deuxième est multipliée par six pour restituer
\(C^*_{\rm deg}\). On utilise enfin
\(\delta_{\rm deg}=9A_{\rm deg}\).
\end{proof}

L'inversion reconstruit les coordonnées du lecteur une fois celui-ci fixé. Elle ne
transforme pas les angles observés en antécédents de \(A\), \(C^*\) ou
\(\Delta_4\).

\subsection{Décomposition géométrique sous inversion de feuillet}
\label{sec:ii09-descomposicion-paridad}

Dans la coordinatisation angulaire, l’inversion de feuillet conserve \(a,d\) et
inverse \(u\). Le vecteur exprimé admet la décomposition
\[
 \boldsymbol\theta_+=a c_A+u c_h+d c_\Delta,
 \qquad
 \boldsymbol\theta_-=a c_A-u c_h+d c_\Delta.
\]
\begin{proposition}[Composantes paire et impaire du transport angulaire]
\label{prop:ii09-paridad-angular}
Les deux vecteurs satisfont
\[
 \frac{\boldsymbol\theta_++\boldsymbol\theta_-}{2}
       =a c_A+d c_\Delta,\qquad
 \frac{\boldsymbol\theta_+-\boldsymbol\theta_-}{2}=u c_h.
\]
La contribution qui change de signe occupe la droite
\(\operatorname{span}(c_h)\) ; la composante conservée occupe
\(\operatorname{span}(c_A,c_\Delta)\).
\end{proposition}
\begin{proof}
La somme annule les termes \(u c_h\), et la différence annule
les deux autres. Les colonnes sont linéairement indépendantes, puisque
\(\det\mathsf M_{\rm CKM}=-3857/6\). Les deux
composantes forment donc une décomposition directe.
\end{proof}

La même inversion agit sur l’évaluation de Barbero--Immirzi :
elle échange \(q_+\) et \(q_-\), et inverse la composante bilinéaire
dans \(A\) et \(C^*\). Par conséquent, la fonctionnelle d’incidence
est impaire. Une observable symétrique dans les deux canaux conserve sa
valeur. Cette différence de parité exprime comment une même
opération sur la structure produit des réponses différentes selon
le lecteur : orientation aréale, information symétrique et vecteur de
mélange conservent des informations différentes du même antécédent.


\subsection{Transport de l'inversion de feuillet}

La conjugaison des coordinatisations angulaires maintient \(A_{\rm deg}\) et
\(\Delta_{\rm deg}\), et inverse \(h_\theta\). Nous définissons
\begin{equation}
 D_\theta=\operatorname{diag}(1,-1,1),\qquad
 c_h=(-5,7,-20)^{\mathsf T},\qquad
 \ell_h=\frac1{3857}(75,176,-150).
 \label{eq:ii-ckm-sheet-data}
\end{equation}
Si \(c_A=(2,0,1)^{\mathsf T}\) et
\(c_\Delta=(28,-25/2,-2/3)^{\mathsf T}\), alors
\(\ell_hc_A=\ell_hc_\Delta=0\) et \(\ell_hc_h=1\).

\begin{theorem}[Réflexion rationnelle et métrique transportée]
L'opérateur
\begin{equation}
 \mathcal R_{\rm CKM}
 =\mathsf M_{\rm CKM}D_\theta\mathsf M_{\rm CKM}^{-1}
 =I_3-2c_h\ell_h
 \label{eq:ii-ckm-reflection}
\end{equation}
satisfait
\begin{equation}
 \mathcal R_{\rm CKM}^2=I_3,\qquad
 \det\mathcal R_{\rm CKM}=-1,\qquad
 \mathcal R_{\rm CKM}\mathsf M_{\rm CKM}
 =\mathsf M_{\rm CKM}D_\theta.
 \label{eq:ii-ckm-intertwining}
\end{equation}
C'est une réflexion orthogonale pour la métrique positive
\begin{equation}
 G_\theta=(\mathsf M_{\rm CKM}^{-1})^{\mathsf T}
                 \mathsf M_{\rm CKM}^{-1},\qquad
 \mathcal R_{\rm CKM}^{\mathsf T}G_\theta\mathcal R_{\rm CKM}=G_\theta.
 \label{eq:ii-ckm-metric}
\end{equation}
\end{theorem}

\begin{proof}
L'opérateur \(P_h=c_h\ell_h\) vérifie \(P_h^2=P_h\), a pour image
\(\mathbb Rc_h\) et pour noyau \(\operatorname{span}(c_A,c_\Delta)\).
Ainsi, \((I_3-2P_h)^2=I_3\), avec valeurs propres \(1,1,-1\).
L'action sur ces trois colonnes est celle de \(D_\theta\) dans la base
originale, ce qui démontre la conjugaison et l'entrelacement.
Pour \(v\ne0\),
\(v^{\mathsf T}G_\theta v=\|\mathsf M_{\rm CKM}^{-1}v\|^2>0\).
La dernière identité résulte de \(D_\theta^{\mathsf T}D_\theta=I_3\).
\end{proof}

Son expression complète, conservée du support précédent, est
\begin{equation}
 \mathcal R_{\rm CKM}=
 \begin{pmatrix}
 4607/3857&1760/3857&-1500/3857\\
 -150/551&199/551&300/551\\
 3000/3857&7040/3857&-2143/3857
 \end{pmatrix}.
 \label{eq:ii-ckm-reflection-matrix}
\end{equation}
La phase \(\delta_{\rm deg}=9A_{\rm deg}\) reste fixe sous cette
réflexion. Le transport de feuillet ne doit donc pas être identifié
à la conjugaison CP des amplitudes, qui change le signe de la phase complexe.
La parité combinatoire d'un chemin, sa conjugaison de feuillet et sa
représentation physique maintiennent les domaines distingués dans
\cite{cpt,mezcla}.

\subsection{Parité du registre et réalisation CP}

Pour conserver le support discret de cette distinction, soit \(\rho\) un
chemin enregistré de longueur \(L\), avec directions
\(d_t\in\{N,E,S,W\}\) et feuillets \(\ell_t\in\{\Sigma,\Pi\}\).
Les états intermédiaires sont ceux produits par le transport TPK.
La réflexion \(P\) agit par \((i,j)\mapsto(i,-j)\pmod9\) et
permute \(E,W\) ; \(C_{\rm hoja}\) permute \(\Sigma,\Pi\) ;
\(T_{\rm reg}\) inverse l'ordre du registre et remplace chaque
direction par son opposée. L'inversion s'effectue sur le registre
enrichi, non sur une projection ayant oublié son prédécesseur.

\begin{proposition}[Invariance de l'action combinatoire de chemin]
Pour \(\lambda\ge0\), la fonctionnelle adimensionnelle
\begin{equation}
 \mathcal A_\lambda(\rho)=
 \sum_{t=2}^{L}\left[1-\delta_{d_t,d_{t-1}}
             +\lambda(1-\delta_{\ell_t,\ell_{t-1}})\right]
 \label{eq:ii-discrete-route-action}
\end{equation}
est invariante sous \(P\), \(C_{\rm hoja}\), \(T_{\rm reg}\) et
leurs compositions.
\end{proposition}

\begin{proof}
Une permutation globale de l'alphabet conserve l'égalité entre voisins.
L'inversion de l'ordre préserve le nombre de changements consécutifs.
Chaque opération conserve donc séparément les changements de direction
et les changements de feuillet sommés par \eqref{eq:ii-discrete-route-action}.
\end{proof}

\(\lambda\) est le paramètre déclaré de cette famille de fonctionnelles ;
\(\mathcal A_\lambda\) ne s’identifie pas à une section dimensionnelle de
Planck. L’invariance du compteur n’implique pas non plus que chaque lecteur
orienté soit pair. Dans une réalisation de mélange \(\mathcal V\), la
compatibilité CP prend la forme précise
\begin{equation}
 \mathcal V(CP_{\rm HMT}\rho)
 =D_u\,\overline{\mathcal V(\rho)}\,D_d^\dagger,
 \label{eq:ii-cp-square}
\end{equation}
avec les changements diagonaux de phase et les espaces de représentation
déclarés. Le compteur de chemin, le transport rationnel de feuillet et la
conjugaison complexe sont trois applications distinctes du support préservé.

\begin{remark}[Compatibilité CP du lecteur physique]
L’équation \eqref{eq:ii-cp-square} est la condition de
compatibilité CP du lecteur physique complet.
Les démonstrations de cette section établissent l’invariance du
registre et les identités du lecteur complexe, mais ne transforment pas
automatiquement \(\mathcal R_{\rm CKM}\) en \(CP_{\rm HMT}\).
Cette réserve concerne ce carré de représentation, non
l’existence du transport interne de chemins déjà construit.
\end{remark}

\subsection{Amplitudes unitaires, phase CP et invariant de Jarlskog}

Notons \(c_{ab}=\cos\vartheta_{ab}\) et
\(s_{ab}=\sin\vartheta_{ab}\). La réalisation complexe du lecteur est
\begingroup
\small
\setlength{\arraycolsep}{3pt}
\begin{equation}
 V_{\rm CKM}=
 \begin{pmatrix}
 c_{12}c_{13}&s_{12}c_{13}&s_{13}e^{-i\delta}\\
 -s_{12}c_{23}-c_{12}s_{23}s_{13}e^{i\delta}&
 c_{12}c_{23}-s_{12}s_{23}s_{13}e^{i\delta}&s_{23}c_{13}\\
 s_{12}s_{23}-c_{12}c_{23}s_{13}e^{i\delta}&
 -c_{12}s_{23}-s_{12}c_{23}s_{13}e^{i\delta}&c_{23}c_{13}
 \end{pmatrix}.
 \label{eq:ii-ckm-unitary}
\end{equation}
\endgroup
La structure complexe utilisée est la réalisation de la structure
orientée déjà construite ; la paramétrisation unitaire représente le résultat
angulaire, elle ne détermine pas ses coefficients.

\begin{proposition}[Unitarité et invariant de phase]
La matrice \eqref{eq:ii-ckm-unitary} est unitaire. Le quartet
\begin{equation}
 J_{\rm CKM}
 =\operatorname{Im}
   (V_{11}V_{22}\overline{V_{12}}\,\overline{V_{21}})
 =c_{12}c_{23}c_{13}^2s_{12}s_{23}s_{13}\sin\delta
 \label{eq:ii-jarlskog}
\end{equation}
est invariant sous \(V\mapsto D_uVD_d\), avec \(D_u,D_d\) diagonales
unitaires, et change de signe sous \(V\mapsto\overline V\).
\end{proposition}

\begin{proof}
La matrice est le produit \(R_{23}U_{13}(\delta)R_{12}\), où
\begin{equation}
\begin{aligned}
 R_{12}&=\begin{pmatrix}c_{12}&s_{12}&0\\-s_{12}&c_{12}&0\\0&0&1\end{pmatrix},
 \\
 U_{13}(\delta)&=
 \begin{pmatrix}c_{13}&0&s_{13}e^{-i\delta}\\0&1&0\\
 -s_{13}e^{i\delta}&0&c_{13}\end{pmatrix},
 \\
 R_{23}&=\begin{pmatrix}1&0&0\\0&c_{23}&s_{23}\\0&-s_{23}&c_{23}\end{pmatrix}.
\end{aligned}
 \label{eq:ii-ckm-factorization}
\end{equation}
Chaque facteur est unitaire par \(c_{ab}^2+s_{ab}^2=1\).
En substituant les quatre entrées du quartet, sa partie imaginaire est
\(c_{12}s_{12}c_{13}^2c_{23}s_{23}s_{13}
(c_{12}^2+s_{12}^2)\sin\delta\), ce qui donne
\eqref{eq:ii-jarlskog}. Les facteurs de phase de chaque ligne et de chaque colonne
se simplifient avec leurs conjugués. La conjugaison complexe inverse la partie
imaginaire.
\end{proof}

Un \(J_{\rm CKM}\) non nul empêche de rendre toute la matrice réelle par
ces changements de phase. Son interprétation physique maintient en outre le secteur
et la représentation correspondants ; ce n'est pas un correcteur scalaire ajouté
aux masses. Si les spectres présentent des dégénérescences, leurs libertés de
changement de base doivent être conservées avant de formuler une conclusion CP.

\subsection{Conjugaison des amplitudes et invariant orienté}
\label{sec:ii09-conjugacion-amplitudes}

La réalisation complexe conserve les trois angles de mélange et
la phase d'interférence comme coordonnées explicites. Écrivons
\(V(\boldsymbol\vartheta,\delta)
=R_{23}U_{13}(\delta)R_{12}\), avec les rotations et la phase
définies dans le développement unitaire. Les arguments des
fonctions trigonométriques sont leurs représentants en radians.

\begin{proposition}[Conjugaison dans la réalisation du mélange]
\label{prop:ii09-conjugacion-unitaria}
Pour les trois angles réels et la phase réelle déclarés,
\[
 V(\boldsymbol\vartheta,-\delta)
       =\overline{V(\boldsymbol\vartheta,\delta)},
 \qquad
 J(\boldsymbol\vartheta,-\delta)
       =-J(\boldsymbol\vartheta,\delta).
\]
L'invariant \(J\) conserve en outre sa valeur sous les changements
diagonaux de phase des bases chargées.
\end{proposition}
\begin{proof}
Les matrices \(R_{12}\) et \(R_{23}\) ont des entrées réelles.
Dans \(U_{13}\), le changement \(\delta\mapsto-\delta\) conjugue
les facteurs exponentiels. La conjugaison du produit donne la
première identité. Dans le quartet
\(V_{11}V_{22}\overline{V_{12}}\overline{V_{21}}\), les
facteurs diagonaux de chaque ligne et colonne s'annulent. La
conjugaison inverse le signe de sa partie imaginaire. On obtient
le résultat de manière équivalente par substitution dans
\[
 J=c_{12}c_{23}c_{13}^{2}s_{12}s_{23}s_{13}\sin\delta.
\]
\end{proof}

Dans le secteur à trois générations, la valeur interne non nulle de
\(J\) exprime l'obstruction à réaliser toutes les amplitudes
par une matrice réelle au moyen de changements diagonaux de phase. C'est
la lecture de violation de charge--parité de la matrice
construite. La lecture des routes conserve, à son tour, l'opération
discrète de conjugaison et de parité ainsi que sa représentation. Leur
compatibilité complète s'exprime par le carré
\[
 \mathcal V(CP_{\rm HMT}\rho)
  =D_u\,\overline{\mathcal V(\rho)}\,D_d^{\dagger},
\]
avec les espaces de représentation et les changements de phase
déclarés. L'identité de conjugaison prouvée ici fixe
exactement l'action que doit transporter ce carré.


\Needspace{14\baselineskip}
\begin{samepage}
\begin{proposition}[Mélange et spectre]
Soient \(B_u=\operatorname{diag}(a_1,a_2,a_3)\) et
\(B_d=\operatorname{diag}(b_1,b_2,b_3)\) les lecteurs hermitiens de masse
déjà produits dans leurs fibres. Le transport
\begin{equation}
 B_d^{(u)}=V_{\rm CKM}B_dV_{\rm CKM}^{\dagger}
 \label{eq:ii-ckm-mass-transport}
\end{equation}
conserve le spectre de \(B_d\) et satisfait
\begin{equation}
 \left|\det[B_u,B_d^{(u)}]\right|
 =2|J_{\rm CKM}|
   \prod_{i<j}|a_i-a_j|\prod_{r<s}|b_r-b_s|.
 \label{eq:ii-ckm-commutator}
\end{equation}
\end{proposition}
\end{samepage}

\begin{proof}
La conjugaison unitaire conserve le polynôme caractéristique. Si
\(B=B_d^{(u)}\), les entrées du commutateur sont
\([B_u,B]_{ij}=(a_i-a_j)B_{ij}\), avec diagonale nulle. Son déterminant est
\[
 2i(a_1-a_2)(a_2-a_3)(a_3-a_1)
     \operatorname{Im}(B_{12}B_{23}B_{31}).
\]
En substituant \(B_{ij}=\sum_r b_rV_{ir}\overline{V_{jr}}\), la partie
imaginaire est un polynôme alterné de degré trois dans les \(b_r\).
Il est divisible par leurs trois différences : si deux \(b_r\) coïncident,
\(B=bI+(b'-b)vv^\dagger\) et le produit hors diagonale est réel.
Le développement d'un coefficient, en utilisant l'orthogonalité des colonnes,
donne le quartet de \eqref{eq:ii-jarlskog}, avec le signe de l'ordre choisi.
Prendre les valeurs absolues élimine cette convention et produit
\eqref{eq:ii-ckm-commutator}, y compris son annulation aux dégénérescences.
\end{proof}

\begin{remark}[Évaluation et comparaison]
Les coordonnées décimales CKM du
tableau~\ref{tab:ii-ckm-comparacion} sont des évaluations ultérieures du
lecteur précédent. La section de comparaison quantitative présente également
les angles et les modules des amplitudes. Ces évaluations ne sont pas
utilisées pour choisir ses colonnes ni pour démontrer
\eqref{eq:ii-ckm-intertwining} ou \eqref{eq:ii-jarlskog}.
Ce bloc n'attribue pas de signification métrologique nouvelle à ces nombres :
une comparaison conjointe demande son édition externe, la définition des
observables et leur covariance. Les écarts marginaux ne la remplacent pas.
\end{remark}

\subsection{Registre leptonique, opérateurs et repères spectraux}

Le secteur leptonique conserve un autre domaine. La construction spectrale de
\eqref{eq:ii-neutrino-compression} réunit quatre profondeurs neutres de dimensions
\(2,4,6,8\), avant la compression tridimensionnelle :
\begin{equation}
 \mathcal H_\nu^{\rm reg}
 =\bigoplus_{j=1}^{4}\mathcal H_{\nu,G_j},\qquad
 \dim\mathcal H_\nu^{\rm reg}=20,\qquad
 M_\nu=\left.(P_\nu\mathcal M_\nu P_\nu)\right|_{\operatorname{Im}P_\nu},
 \quad\operatorname{rank}P_\nu=3.
 \label{eq:ii-neutrino-compression}
\end{equation}
Les quatre profondeurs ne sont ni quatre saveurs ni quatre masses.
\(P_\nu\) est le projecteur neutre hérité du lecteur de registre, et
\(M_\ell=\left.(P_\ell\mathcal M_\ell P_\ell)\right|_{\operatorname{Im}P_\ell}\),
avec \(\operatorname{rank}P_\ell=3\), est l'opérateur chargé comprimé.
L'échelle de ces opérateurs provient de la branche action--horloge--masse ;
aucun angle PMNS ne la détermine.

\begin{definition}[Transport des repères leptoniques]
Dans la réalisation tridimensionnelle commune déclarée, soient
\(F_\ell,F_\nu:\mathbb C^3\to\mathcal H_3\) des repères spectraux
orthonormaux complets des opérateurs chargé et neutre. On fixe
l'identification avec \(\mathcal H_3\) avant de prendre leur produit. Le lecteur est
\begin{equation}
 U_{\rm PMNS}=F_\ell^\dagger F_\nu:
 \mathbb C^3_\nu\longrightarrow\mathbb C^3_\ell.
 \label{eq:ii-pmns-frames}
\end{equation}
Les repères et l'identification proviennent de la représentation et du
registre ; ils ne sont pas choisis à partir d'une matrice expérimentale que l'on souhaite
reproduire.
\end{definition}

\begin{theorem}[Unitarité, dégénérescences et naturalité]
Sous l'identification précédente, \(U_{\rm PMNS}\) est unitaire.
Si les spectres sont simples, leurs changements de repère sont des phases diagonales
et des permutations ; l'ordre spectral étant fixé, les phases demeurent.
S'il existe des multiplicités \(m_{\ell,\lambda}\) et \(m_{\nu,\eta}\),
l'objet indépendant de ces bases est la classe double
\begin{equation}
 \left(\prod_\lambda U(m_{\ell,\lambda})\right)
 \backslash U(3) /
 \left(\prod_\eta U(m_{\nu,\eta})\right).
 \label{eq:ii-pmns-double-class}
\end{equation}
Un raffinement unitaire commun \(W\) qui entrelace les deux opérateurs
conserve le transport lorsque les repères se prolongent selon
\(F'_\ell=WF_\ell\), \(F'_\nu=WF_\nu\).
\end{theorem}

\begin{proof}
La complétude dans le même espace donne
\(F_\ell F_\ell^\dagger=F_\nu F_\nu^\dagger=I_{\mathcal H_3}\).
Par conséquent,
\(U_{\rm PMNS}^\dagger U_{\rm PMNS}
=F_\nu^\dagger F_\ell F_\ell^\dagger F_\nu=I_3\), et de même
dans l'autre ordre. Les changements de base préservant un opérateur hermitien
agissent par blocs dans ses espaces propres. En substituant
\(F_\ell D_\ell,F_\nu D_\nu\), on obtient
\(U\mapsto D_\ell^\dagger UD_\nu\), ce qui détermine la classe double.
Enfin, \((WF_\ell)^\dagger WF_\nu=F_\ell^\dagger F_\nu\).
\end{proof}

La condition d'espace commun est effective. Pour deux repères isométriques dans
un espace ambiant plus grand, l'orthonormalité séparée ne produit qu'une
compression ; elle n'autorise pas à simplifier \(F_\ell F_\ell^\dagger\) comme
s'il s'agissait de l'identité de l'espace ambiant.

\subsection{Noyau de registre et suffisance du rang}

Les lecteurs \(\Xi_r\) conservent les composantes de retenue arithmétique, de feuillet,
d'orientation, de lacune et d'holonomie du registre. Sur l'ensemble fini
des chemins admis, on forme
\begin{equation}
 K_{\rm reg}(c,d)=
 \sum_{\substack{r\in\{\mathrm{car},\mathrm{hoj},\mathrm{ori},
                           \mathrm{vac},\mathrm{hol}\}\\ d_r>0}}
 \frac{\langle\Xi_r(c),\Xi_r(d)\rangle}{d_r},
 \qquad
 d_r=\sum_x\|\Xi_r(x)\|^2.
 \label{eq:ii-reg-kernel}
\end{equation}
La normalisation provient des normes du registre lui-même. Les
composantes de norme nulle sont omises ; elles ne sont pas remplacées par des termes tirés
d'une matrice cible. Chaque \(\Xi_r\) et chaque classe d'agrégation doit
conserver son rattachement au chemin et sa compatibilité avec le raffinement.

\begin{proposition}[Positivité et rang du registre]
La matrice \(K_{\rm reg}\) est positive semi-définie. Si ses
vecteurs de caractéristiques sont agrégés en trois secteurs, la matrice de Gram
\(G_\nu\) résultante est inversible exactement lorsque ces trois vecteurs
agrégés sont linéairement indépendants.
\end{proposition}

\begin{proof}
On définit
\(v_c=\bigoplus_{r:d_r>0}d_r^{-1/2}\Xi_r(c)\).
Alors \(K_{\rm reg}(c,d)=\langle v_c,v_d\rangle\), et pour
tous coefficients \(z_c\),
\(\sum_{c,d}\overline z_cK_{\rm reg}(c,d)z_d
=\|\sum_c z_cv_c\|^2\ge0\).
La même identité vaut pour les trois vecteurs agrégés.
Le noyau de leur matrice de Gram est exactement constitué de leurs relations
linéaires, d'où le critère d'inversibilité.
\end{proof}

On distingue cette matrice de Gram d’un secteur de la matrice \emph{croisée} entre les
classes chargées \(L_a\) et neutres \(N_b\), exprimée dans \cite{mezcla} :
\begin{equation}
 G_{ab}=\frac1{\sqrt{|L_a||N_b|}}
       \sum_{c\in L_a}\sum_{d\in N_b}K_{\rm reg}(c,d).
 \label{eq:ii-cross-gram}
\end{equation}
Les classes non vides proviennent du classificateur interne. La matrice croisée
n'est pas nécessairement hermitienne, contrairement à une matrice de Gram construite avec
les mêmes trois vecteurs des deux côtés.

\begin{proposition}[Facteur polaire et sa portée]
Si \(G\) est inversible,
\begin{equation}
 U_G=G(G^\dagger G)^{-1/2}
 \label{eq:ii-polar}
\end{equation}
est unitaire. Si \(G\) est de plus hermitienne définie positive, alors
\(U_G=I_3\). Le rang d'une matrice de Gram positive certifie donc
la suffisance des observables, mais ne détermine pas à lui seul un
mélange non trivial entre les repères chargé et neutre.
\end{proposition}

\begin{proof}
\(G^\dagger G\) est définie positive ; sa racine positive est inversible et
\[
 U_G^\dagger U_G=(G^\dagger G)^{-1/2}(G^\dagger G)
 (G^\dagger G)^{-1/2}=I_3.
\]
Par dimension finie, on a également
\(U_GU_G^\dagger=I_3\). Si \(G>0\), la racine positive de
\(G^\dagger G=G^2\) est \(G\), et \(U_G=I_3\).
\end{proof}

\begin{remark}[Factorisation polaire et repères spectraux]
Les équations \eqref{eq:ii-reg-kernel} et \eqref{eq:ii-cross-gram}
conservent respectivement le critère de rang du registre et
la matrice croisée des secteurs chargé et neutre. Cet exposé
sépare leurs fonctions pour ne pas transformer une matrice de Gram positive en sélecteur
artificiel d’angles. Identifier \(U_G\) à
\eqref{eq:ii-pmns-frames} nécessite la compatibilité du noyau croisé
avec les repères spectraux déclarés. L’identité d’un facteur polaire
ou la seule unitarité ne remplacent pas cette identification. Aucune
nouvelle matrice numérique PMNS n’est exprimée ici.
\end{remark}

\subsection{Contrôle de la phase de rang un}

Soit \(b_K\) le témoin normalisé d'incidence de la source et
\(Q_K=|b_K\rangle\langle b_K|\). L'opérateur
\begin{equation}
 \Phi_K=I+(e^{i\pi_{\rm HMT}/6}-1)Q_K
 \label{eq:ii-rank-one-phase}
\end{equation}
est unitaire dans son espace ambiant. Pour des repères isométriques, on considère
\(U_\Phi=F_\ell^\dagger\Phi_KF_\nu\).

\begin{proposition}[Critère exact d'une compression unitaire]
L'opérateur \(U_\Phi\) est une contraction. Il est unitaire exactement
lorsque
\begin{equation}
 \Phi_K(\operatorname{Im}F_\nu)=\operatorname{Im}F_\ell.
 \label{eq:ii-phase-compatibility}
\end{equation}
\end{proposition}

\begin{proof}
Soient \(P_\ell=F_\ell F_\ell^\dagger\) et
\(P_\nu=F_\nu F_\nu^\dagger\). Alors
\(U_\Phi^\dagger U_\Phi
=F_\nu^\dagger\Phi_K^\dagger P_\ell\Phi_KF_\nu\le I_3\).
L'égalité équivaut à
\((I-P_\ell)\Phi_KF_\nu=0\), c'est-à-dire à l'inclusion de
l'image de gauche de \eqref{eq:ii-phase-compatibility} dans celle de droite.
Les deux ont la dimension trois, de sorte que l'inclusion est une égalité.
\end{proof}

La source \cite{mezcla} conserve cet ansatz minimal comme contrôle
négatif d'identification physique. Son unitarité ambiante et même la
satisfaction de \eqref{eq:ii-phase-compatibility} ne remplacent pas le
registre complet ni les repères de \eqref{eq:ii-pmns-frames}. Ce
contrôle n'est pas éliminé lors de la réunion de la propagation angulaire dans l'article II.

\begingroup\fontsize{10}{11.7}\selectfont
\setlength{\parskip}{2pt}\setlength{\abovedisplayskip}{5pt}\setlength{\belowdisplayskip}{5pt}
\subsection{Phase leptonique et critère Dirac--Majorana}

Les changements de phase des bases doivent être distingués des données
d'autoconjugaison du secteur neutre. Pour tout transport unitaire
leptonique, on peut former le quartet
\begin{equation}
 J_\ell=\operatorname{Im}
 \bigl(U_{e1}U_{\mu2}\overline{U_{e2}}\,
                        \overline{U_{\mu1}}\bigr).
 \label{eq:ii-leptonic-quartet}
\end{equation}
La simplification des phases démontrée pour \eqref{eq:ii-jarlskog}
s'applique également ici. Le quartet ne fixe pas toutes les données de la
réalisation neutre et ne détermine pas à lui seul son type Dirac ou Majorana.

La réalisation de Majorana est spécifiée par
l'autoconjugaison antiunitaire \(\mathcal C_\nu\), compatible avec la masse et la
dynamique :
\begin{equation}
 \mathcal C_\nu^2=I,\qquad
 \mathcal C_\nu M_\nu=M_\nu\mathcal C_\nu,\qquad
 \mathcal C_\nu H_\nu=H_\nu\mathcal C_\nu,
 \label{eq:ii-majorana-compatibility}
\end{equation}
avec domaines préservés. La représentation se réalise dans la forme réelle
fixe correspondante. L'alternative Dirac conserve en outre une charge
\(Q\) distinguant le mode de son conjugué :
\begin{equation}
 \mathcal C_\nu Q\mathcal C_\nu^{-1}=-Q,\qquad Q\psi\ne0.
 \label{eq:ii-dirac-charge}
\end{equation}
Ces énoncés conservent leurs données de représentation. La neutralité,
l'involution des mots dodécaphasiques et un secteur topologique d'Ising
ne décident pas séparément de l'alternative.

\begin{remark}[Liberté des bases et phases de Majorana]
La classe double \eqref{eq:ii-pmns-double-class} décrit la liberté des
bases du couple hermitien. Si l'on fixe en outre une forme réelle de Majorana,
les changements de phase admissibles doivent préserver cette donnée supplémentaire ;
on ne peut pas y transférer sans autre examen toutes les libertés d'une représentation
de Dirac. Le présent bloc conserve le critère d'autoconjugaison et
n'attribue de valeurs ni aux phases de Majorana ni aux masses neutriniques. Ces valeurs
ne s'obtiennent ni de la neutralité ni d'une polarisation de rang un.
\end{remark}

\Needspace{8\baselineskip}
\subsection{Clôture du transport angulaire}

La construction du secteur des quarks détermine d’abord les coordonnées angulaires, puis
les amplitudes complexes et, enfin, les invariants de phase et le
transport d’opérateurs. La construction leptonique détermine des opérateurs de
registre et des repères spectraux avant le changement de base. Les deux chemins
conservent leur résidence dans l’état enrichi commun et dans la structure
discrète du continu ; aucun n’introduit de tableau expérimental comme
sélecteur des coefficients ou des états.

Les contrôles de cette section sont locaux et explicites : inversion et
entrelacement rationnel CKM ; conservation spectrale par conjugaison ;
complétude et espace commun des repères PMNS ; rang du registre ;
compatibilité de la matrice croisée ; et préservation des données
d'autoconjugaison. La propagation angulaire s'intègre ainsi à la théorie de
l'action, de l'aire et de l'information sans faire du mélange une seconde loi des
masses ni transformer le lecteur de Barbero en sa cause.

\par\endgroup
